% ══════════════════════════════════════════════════════════════════════════════
\chapter{Anti-dictionnaire}
\label{chap:td2}
% ══════════════════════════════════════════════════════════════════════════════

À partir du fichier \texttt{corpus.xml} produit en TD1, la deuxième étape vise à
identifier et supprimer les termes non porteurs de sens avant toute indexation. Cette
liste de mots exclus, appelée \textbf{anti-dictionnaire} ou \emph{stop list}, est au
cœur de la gestion du vocabulaire en RI (cours, chapitre~2). La méthode retenue est
statistique et s'appuie sur le coefficient TF-IDF, qui quantifie formellement le
caractère discriminant de chaque terme dans le corpus.

\section{Fondements théoriques : TF-IDF et loi de Zipf}

La loi de Zipf décrit la distribution des fréquences de mots dans tout grand corpus :
un petit nombre de termes très fréquents représente la majorité des occurrences mais
véhicule peu d'information thématique. La mesure formelle retenue est le coefficient
\textbf{IDF} (\emph{Inverse Document Frequency}) :
\begin{equation}
  \operatorname{idf}(t) = \log_{10}\!\left(\frac{N}{\operatorname{df}(t)}\right)
  \label{eq:idf}
\end{equation}
où $N = 326$ est le nombre d'articles et $\operatorname{df}(t)$ le nombre d'articles
contenant au moins une occurrence du terme $t$. Un token ubiquitaire (présent dans
tous les articles) obtient $\operatorname{idf}(t)=0$, l'excluant naturellement de tout
calcul de pertinence. La pondération finale est le produit :
\begin{equation}
  \operatorname{tf\text{-}idf}(t,d)
    = \bigl(1 + \log_{10} \operatorname{count}(t,d)\bigr)
      \times \operatorname{idf}(t)
  \label{eq:tfidf}
\end{equation}
La TF logarithmique (formule~\ref{eq:tfidf}) comprime l'influence des termes très
répétés, conformément au cours (§5.4) : 10 occurrences donnent $\operatorname{tf}=2{,}0$
là où une TF brute aurait donné 10.

\section{Pipeline de calcul et seuillage}

Le pipeline s'exécute en six étapes séquentielles. La tokenisation repose sur la
regex \texttt{[a-zA-ZÀ-ÿ]+} et inclut une normalisation préalable des élisions
françaises, essentielle pour éviter les tokens parasites mono-caractères :

\begin{lstlisting}[caption={Normalisation des élisions avant tokenisation},
  label={lst:elisions}]
_ELISION_RE = re.compile(r"\b(qu|[ldjmtscn])['']", re.IGNORECASE)

def normalize_elisions(text: str) -> str:
    """Supprime les préfixes clitiques : l'innovation -> innovation."""
    return _ELISION_RE.sub("", text)
\end{lstlisting}

Sans cette normalisation, \texttt{l'innovation} produirait les tokens \texttt{l} et
\texttt{innovation}, gonflant artificiellement la fréquence de résidus mono-caractères
(\texttt{l}, \texttt{d}, \texttt{s}…) dans l'IDF. Le seuillage applique deux critères
inspirés de la loi de Zipf : un \textbf{seuil bas} $\theta=0{,}5$ pour exclure les
termes trop communs ($\operatorname{df}(t) \geq 103$ articles, soit 31\,\% du corpus),
et un \textbf{seuil haut} désactivé ($+\infty$) afin de préserver les hapaxes qui
sont souvent les meilleurs discriminants pour les requêtes spécialisées.

\section{Résultats et analyse critique}

La tokenisation du corpus produit 163\,842 occurrences et 14\,177 formes distinctes,
soit environ 503 tokens par article. Le seuil $\theta=0{,}5$ identifie \textbf{88 mots
vides}, exclusivement des articles, prépositions, pronoms et auxiliaires du français.

\begin{figure}[ht]
  \centering
  \includegraphics[width=0.82\textwidth]{../outputs/plot_idf_lemmes.png}
  \caption{Distribution des scores IDF sur les lemmes du corpus ADIT.
    La concentration à $\operatorname{idf}=2{,}51$ illustre la prédominance des termes
    techniques à faible fréquence documentaire (médiane = maximum).}
  \label{fig:idf_distribution}
\end{figure}

La figure~\ref{fig:idf_distribution} révèle une distribution bimodale caractéristique
d'un corpus de veille spécialisé : un petit groupe de tokens ubiquitaires (IDF $\approx 0$)
et une masse de termes techniques n'apparaissant que dans un ou deux articles
(IDF $\approx 2{,}51$). Plus de 68\,\% du vocabulaire n'apparaît que dans 1 à 2
articles — ces quasi-hapaxes constituent les descripteurs les plus discriminants du
corpus et justifient la désactivation du seuil haut. Le top des scores TF-IDF confirme
la cohérence sémantique : \texttt{hoffmann} (5,67), \texttt{sérotonine} (5,47),
\texttt{envisat} (5,39) sont tous des termes très spécifiques à un unique article.
La couverture de tests atteint 95\,\% sur 165 tests.
